\cleardoublepage

\chapter*{Course Map: Follow the Flow of Information}
\addcontentsline{toc}{chapter}{Course Map: Follow the Flow of Information}
\markboth{Course Map}{}

ERP becomes much easier once you stop treating each module as an isolated
topic. The course follows the same logic as a real business: understand the
organisation, capture demand, plan and obtain resources, fulfil the order,
record the financial consequences, and improve the process.

\begin{erpfig}[Course structure]{The five stages of the course. Each stage
consumes what the previous one produced, which is also the order in which
information moves through a real company.}{fig:course-stages}
\begin{tikzpicture}[
  stage/.style={draw=UTPBlue,fill=UTPPaleBlue,rounded corners=2mm,text width=11.9cm,
                minimum height=1.20cm,align=left,inner xsep=4mm,inner ysep=2.3mm,
                font=\scriptsize\color{UTPNavy}},
  node distance=0.52cm
]
\node[stage] (a) {{\normalsize\bfseries\color{UTPBlue}Stage 1: ERP foundations}\hfill{\footnotesize Chapters 1--2}\par
\smallskip Functional areas, processes, silos, shared databases, ERP evolution, benefits, challenges.};
\node[stage,below=of a] (b) {{\normalsize\bfseries\color{UTPBlue}Stage 2: Customer and sales}\hfill{\footnotesize Chapters 3--4}\par
\smallskip CRM, quotations, sales orders, availability, delivery, invoicing, payment.};
\node[stage,below=of b] (c) {{\normalsize\bfseries\color{UTPBlue}Stage 3: Plan and supply}\hfill{\footnotesize Chapters 5--7}\par
\smallskip Production approaches, forecasting, S\&OP, demand management, BOMs, MRP, purchasing, inventory.};
\node[stage,below=of c] (d) {{\normalsize\bfseries\color{UTPBlue}Stage 4: Finance, people, improvement}\hfill{\footnotesize Chapters 8--10}\par
\smallskip Accounting integration, HR processes, flowcharts, swimlanes, change management, go-live.};
\node[stage,below=of d,draw=UTPGreen,fill=UTPPaleTeal] (e) {{\normalsize\bfseries\color{UTPGreen}Stage 5: Odoo practical integration}\hfill{\footnotesize Labs 1--8 + Capstone}\par
\smallskip Build West End Bicycles in Odoo 19 Community and trace one story from setup through purchasing, manufacturing, CRM, sales, delivery, invoicing, and reporting.};
\draw[erpflow] (a) -- (b); \draw[erpflow] (b) -- (c);
\draw[erpflow] (c) -- (d); \draw[erpflow] (d) -- (e);
\node[erpcap,rotate=90,anchor=south] at (-6.55,-3.2) {information moves this way};
\end{tikzpicture}
\end{erpfig}

\section*{The question behind every chapter}

One business event --- a customer confirming an order --- has to reach four
different desks, each of which needs a different part of it.

\begin{erpfig}[One event, four readings]{The recurring pattern of the course. A
single business event is written once and read four ways. Chapters 1--10 work
through the areas clockwise; the labs rebuild the whole loop in
Odoo.}{fig:one-event-four-readings}
\begin{tikzpicture}[scale=0.98,transform shape]
\node[erpdb,minimum width=3.1cm,minimum height=1.5cm] (db) at (0,0) {Shared\\record};
\node[erpstep,minimum width=3.0cm] (ms) at (-3.9,1.55) {Marketing \& Sales\\[-0.2em]{\tiny\mdseries can we promise it?}};
\node[erpstep,minimum width=3.0cm] (scm) at (3.9,1.55) {Supply Chain\\[-0.2em]{\tiny\mdseries can we make it?}};
\node[erpstep,minimum width=3.0cm] (af) at (3.9,-1.55) {Accounting \& Finance\\[-0.2em]{\tiny\mdseries what is owed?}};
\node[erpstep,minimum width=3.0cm] (hr) at (-3.9,-1.55) {Human Resources\\[-0.2em]{\tiny\mdseries who does the work?}};
\draw[erpflow] (ms) -- (db); \draw[erpflow] (scm) -- (db);
\draw[erpflow] (db) -- (af); \draw[erpflow] (hr) -- (db);
\node[erpchip] at (0,2.05) {Ch.\,1--4};
\node[erpchip] at (0,-2.05) {Ch.\,8--9};
\end{tikzpicture}
\end{erpfig}

